% ============================================================
% Plantilla: Artículo con el sistema Harvard (autor-año)
%   \parencite{clave} -> (Autor, año)      \textcite{clave} -> Autor (año)
% Recompilar ejecuta biber automáticamente.
% ============================================================
\documentclass[12pt,a4paper]{article}

\usepackage[utf8]{inputenc}
\usepackage[T1]{fontenc}
\usepackage{lmodern}
\usepackage[spanish,es-tabla]{babel}
\usepackage[margin=2.5cm]{geometry}
\usepackage{setspace}
\usepackage{csquotes}
\usepackage[style=authoryear,backend=biber,maxcitenames=2,uniquename=false,dashed=false]{biblatex}
\usepackage[hidelinks]{hyperref}

\addbibresource{referencias.bib}
\onehalfspacing

% Harvard: coma entre autor y año, "y" entre autores
\renewcommand*{\nameyeardelim}{\addcomma\space}

\title{Título del artículo}
\author{Nombre del Autor\\\small Universidad}
\date{\today}

\begin{document}

\maketitle

\begin{abstract}
Resume el objetivo, el método y las conclusiones del artículo.
\end{abstract}

\section{Introducción}
El sistema Harvard cita al autor y el año de la obra \parencite{lamport1994}. Cuando el
autor forma parte de la oración, se escribe \textcite{knuth1984}.

\section{Revisión de la literatura}
Varios autores coinciden en esta idea \parencite{goossens2004,lamport1994}.

\section{Conclusiones}
Resume los hallazgos y propone líneas futuras.

\printbibliography[title=Referencias]

\end{document}
